\documentclass[12pt]{article}
\usepackage[margin=1in]{geometry}
\usepackage{amsmath,amssymb}
\usepackage{booktabs}
\usepackage[colorlinks=true,linkcolor=blue]{hyperref}

% Same macros as main.tex
\newcommand{\GB}{\mathrm{GB}}
\newcommand{\eff}{\mathrm{eff}}
\newcommand{\intf}{\mathrm{int}}
\newcommand{\Ox}{\mathrm{O}}
% Local shorthands
\newcommand{\D}{D_{\Ox}^{\eff}}
\newcommand{\Cb}{\bar{C}_{\Ox}}
\newcommand{\Cc}{C_{\Ox}^{\mathrm{c}}}
\newcommand{\Ci}{C_{\Ox}^{\intf}}
\newcommand{\Jr}{J_{\mathrm{r}}}
\newcommand{\dd}{\mathrm{d}}

\title{Cross-section-averaged oxygen transport along an oxidized grain boundary:\\
lateral closure from the separable solution of the Laplace equation}


\begin{document}
\maketitle

\section*{Purpose and result}

The GB oxidation model reduces oxygen transport to one dimension
along the boundary. This note treats the entire oxide cross-section as the
transport channel and derives, from the two-dimensional diffusion equation,
\begin{enumerate}
  \item[(i)] the one-dimensional balance for the cross-section-averaged oxygen
        concentration $\Cb$;
  \item[(ii)] a closure that gives the interfacial flux $\Jr$ directly from $\Cb$, the
        oxide half-width $L$, the diffusivity $\D$ and the interfacial reaction law
        $\Jr(\Ci)$.
\end{enumerate}
The closure is obtained from the separable solution of the Laplace equation in the
oxide cross-section. It reduces to one scalar equation for a dimensionless
eigenvalue $\mu\in(0,\pi/2)$,
\begin{equation}
\boxed{\;
\frac{\D\,\mu^{2}\,\Cb}{L}
 =\sum_k \Jr^{k}\!\bigl(\Cb\,\mu\cot\mu\bigr)
\;}
\label{eq:mu-eq}
\end{equation}
after which
\begin{equation}
\boxed{\;
\Jr=\frac{\D\,\mu^{2}\,\Cb}{L},
\qquad
\Ci=\Cb\,\mu\cot\mu
\;}
\label{eq:result}
\end{equation}
Equation~\eqref{eq:mu-eq} has a unique root and is solved pointwise by Newton's method,
exactly as same as the present quasi-steady closure solved for $\Ci$.

\section{Geometry and notation}
\label{sec:notation}

\begin{itemize}
  \item $y$: coordinate along the grain boundary, $y=0$ at the surface.
  \item $x$: coordinate normal to the boundary, $x=0$ on the boundary mid-plane
        (plane of symmetry). Only the half-domain $0\le x\le L$ is considered.
  \item $L(y,t)$: distance from the mid-plane to the oxide/metal interface. After the original boundary channel is filled.
  \item $C_{\Ox}(x,y,t)$: concentration of mobile (dissolved) oxygen in the oxide.
  \item $\D(y,t)$: effective oxygen diffusivity of the locally mixed oxide,
        Eq.~(\texttt{eq:Deff}) of \texttt{main.tex}. It is uniform across the section,
        because the phases are mixed locally rather than layered.
  \item Cross-section average, mid-plane value and interface value:
        \begin{equation}
        \Cb(y,t)=\frac{1}{L}\int_{0}^{L} C_{\Ox}\,\dd x,\qquad
        \Cc=C_{\Ox}(0,y,t),\qquad
        \Ci=C_{\Ox}(L,y,t).
        \end{equation}
  \item Interfacial oxygen flux (per unit boundary area, positive toward the metal):
        \begin{equation}
        \Jr(y,t)=-\D\,\frac{\partial C_{\Ox}}{\partial x}\bigg|_{x=L}.
        \label{eq:Jr-def}
        \end{equation}
  \item Interfacial reaction law: the total consumption
        $\Jr(\Ci)=\sum_k\Jr^{k}(\Ci)$ of
        Eqs.~(\texttt{eq:rate-cr2o3})--(\texttt{eq:rate-total}) of \texttt{main.tex}.
        It is a nondecreasing function of $\Ci$ at fixed metal concentrations. Where
        a linearized form is used below, $k\equiv\dd\Jr/\dd\Ci$.
\end{itemize}

\section{Assumptions}
\label{sec:assumptions}

\begin{enumerate}
  \item[A1.] Two-dimensional Fickian diffusion of oxygen in the oxide with a single,
        isotropic $\D$ that does not vary across the section.
  \item[A2.] Symmetry at the mid-plane: $\partial C_{\Ox}/\partial x=0$ at $x=0$.
  \item[A3.] Slow variation along the boundary of the half-width, the diffusivity and
        the metal concentrations that enter the reaction law:
        $|\partial L/\partial y|\ll1$, and $\D$, $\Jr(\cdot)$ change little over a
        distance $L$.
  \item[A4.] Lateral quasi-steady state: the lateral profile adjusts on the time scale
        $L^2/\D$, which is short compared with the time scales on which $\Cb$ and $L$
        change.
  \item[A5.] Locally, the lateral profile is the fundamental separable mode of the
        Laplace equation (Sec.~\ref{sec:mode}). Higher modes decay over a distance
        of order $0.3L$ or less and are neglected (Sec.~\ref{sec:validity}).
\end{enumerate}
No assumption is made on whether the interface is reaction- or diffusion-limited, and
no assumption is made that $C_{\Ox}$ is uniform across the section.

\section{Cross-section integration of the 2D balance}
\label{sec:integration}

The two-dimensional conservation equation in the oxide is
\begin{equation}
\frac{\partial C_{\Ox}}{\partial t}
 =\frac{\partial}{\partial y}\!\left(\D\,\frac{\partial C_{\Ox}}{\partial y}\right)
 +\D\,\frac{\partial^2 C_{\Ox}}{\partial x^2}.
\label{eq:2d}
\end{equation}
Integrate each term over $0\le x\le L(y,t)$, using the Leibniz rule where the limit
depends on $y$ or $t$.

\paragraph{Lateral term.} With A2 and the definition \eqref{eq:Jr-def},
\begin{equation}
\int_{0}^{L}\D\,\frac{\partial^2 C_{\Ox}}{\partial x^2}\,\dd x
 =\D\,\frac{\partial C_{\Ox}}{\partial x}\bigg|_{0}^{L}
 =-\Jr .
\label{eq:int-lateral}
\end{equation}

\paragraph{Time term.} Because the upper limit moves, the Leibniz rule gives the
identity
\begin{equation}
\frac{\partial}{\partial t}\int_{0}^{L}C_{\Ox}\,\dd x
 =\int_{0}^{L}\frac{\partial C_{\Ox}}{\partial t}\,\dd x
  +C_{\Ox}(L)\,\frac{\partial L}{\partial t},
\qquad\text{i.e.}\qquad
\int_{0}^{L}\frac{\partial C_{\Ox}}{\partial t}\,\dd x
 =\frac{\partial}{\partial t}\bigl(L\,\Cb\bigr)-\Ci\,\frac{\partial L}{\partial t}.
\label{eq:int-time}
\end{equation}
The last term is the dissolved oxygen, at concentration $\Ci$, contained in the
oxide formed during $\dd t$.

\paragraph{Along-boundary term.} Applying the Leibniz rule twice,
\begin{align}
\int_{0}^{L}\frac{\partial}{\partial y}\!\left(\D\frac{\partial C_{\Ox}}{\partial y}\right)\dd x
 &=\frac{\partial}{\partial y}\int_{0}^{L}\D\frac{\partial C_{\Ox}}{\partial y}\,\dd x
   -\D\,\frac{\partial C_{\Ox}}{\partial y}\bigg|_{x=L}\frac{\partial L}{\partial y},
\\
\int_{0}^{L}\D\frac{\partial C_{\Ox}}{\partial y}\,\dd x
 &=\D\left[\frac{\partial}{\partial y}\bigl(L\,\Cb\bigr)-\Ci\,\frac{\partial L}{\partial y}\right]
  =\D\,L\,\frac{\partial\Cb}{\partial y}
   +\D\,(\Cb-\Ci)\,\frac{\partial L}{\partial y}.
\end{align}
Hence
\begin{equation}
\int_{0}^{L}\frac{\partial}{\partial y}\!\left(\D\frac{\partial C_{\Ox}}{\partial y}\right)\dd x
 =\frac{\partial}{\partial y}\!\left(\D\,L\,\frac{\partial\Cb}{\partial y}\right)
  +\underbrace{\frac{\partial}{\partial y}\!\left[\D(\Cb-\Ci)\frac{\partial L}{\partial y}\right]
   -\D\frac{\partial C_{\Ox}}{\partial y}\bigg|_{x=L}\frac{\partial L}{\partial y}}_{O(\partial L/\partial y)} .
\label{eq:int-y}
\end{equation}

\paragraph{Result.} Collecting \eqref{eq:int-lateral}--\eqref{eq:int-y},
\begin{equation}
\frac{\partial}{\partial t}\bigl(L\,\Cb\bigr)
 =\frac{\partial}{\partial y}\!\left(\D\,L\,\frac{\partial\Cb}{\partial y}\right)
  -\Jr
  +\Ci\,\frac{\partial L}{\partial t}
  +O\!\left(\frac{\partial L}{\partial y}\right).
\label{eq:1d-exact}
\end{equation}
Equation~\eqref{eq:1d-exact} is exact. Under A3 the last group is dropped.

At the moving interface the diffusive flux $\Jr$ of \eqref{eq:Jr-def} supplies both
the reaction and the dissolved oxygen in the fresh oxide (Stefan condition, dissolved
oxygen in the metal neglected):
\begin{equation}
\Jr=\sum_k\Jr^{k}(\Ci)+\Ci\,\frac{\partial L}{\partial t}.
\label{eq:stefan}
\end{equation}
Substituting \eqref{eq:stefan} into \eqref{eq:1d-exact}, the terms
$\Ci\,\partial L/\partial t$ cancel and the sink is the reaction consumption alone.
The working one-dimensional equation is
\begin{equation}
\boxed{\;
\frac{\partial}{\partial t}\bigl(L\,\Cb\bigr)
 =\frac{\partial}{\partial y}\!\left(\D\,L\,\frac{\partial\Cb}{\partial y}\right)
  -\sum_k\Jr^{k}
\;}
\label{eq:1d-working}
\end{equation}
With $\partial L/\partial t=\Omega\sum_k\Jr^k$ and $\Omega=1/C_{\Ox}^{\mathrm{oxide}}$
the oxide volume formed per unit oxygen consumed, the ratio of the Stefan term in
\eqref{eq:stefan} to the reaction term is $\Ci/C_{\Ox}^{\mathrm{oxide}}\ll1$. The
diffusive flux $\Jr$ and the reaction consumption $\sum_k\Jr^k$ therefore differ by a
negligible amount, and from here on the single symbol $\Jr$ is used for both; in
particular the closure of Sec.~\ref{sec:mu} equates $\Jr$ of \eqref{eq:Jr-def} to
$\sum_k\Jr^k$.
Two remarks:
\begin{itemize}
  \item The along-boundary conductance is $\D L$ with the \emph{same} $\D$ as in the
        lateral direction. This follows from A1 alone; no assumption on the lateral
        profile was used.
  \item $L$ sits inside the time derivative. Expanding,
        $L\,\partial\Cb/\partial t=\partial_y(\D L\,\partial_y\Cb)-\Jr-\Cb\,\partial L/\partial t$;
        the last term is the dilution of dissolved oxygen as the section widens.
\end{itemize}
Equation~\eqref{eq:1d-working} contains two unknowns per position, $\Cb$ and $\Jr$.
The relation between them is derived next.

\section{Lateral problem: separation of variables}
\label{sec:mode}

\subsection{Local problem}

Under A3 and A4, in the neighbourhood of a given position along the boundary, the
oxide occupies a half-strip $0\le x\le L$ with $L$, $\D$ and the reaction law treated
as constants, and the concentration satisfies the steady equation
\begin{equation}
\frac{\partial^2 C_{\Ox}}{\partial x^2}+\frac{\partial^2 C_{\Ox}}{\partial y^2}=0,
\qquad
\frac{\partial C_{\Ox}}{\partial x}\bigg|_{x=0}=0,
\qquad
-\D\,\frac{\partial C_{\Ox}}{\partial x}\bigg|_{x=L}=\Jr\bigl(C_{\Ox}(L,y)\bigr).
\label{eq:laplace}
\end{equation}

\subsection{Separable solutions}

Seek $C_{\Ox}=X(x)\,Y(y)$. Substituting into the Laplace equation and dividing by
$XY$,
\begin{equation}
\frac{X''}{X}=-\frac{Y''}{Y}=-\left(\frac{\mu}{L}\right)^2,
\end{equation}
where the separation constant is written as $-(\mu/L)^2$ with $\mu>0$, because the
$x$-problem must be oscillatory (it has a homogeneous Neumann condition at $x=0$ and a
sink at $x=L$) and the $y$-dependence must decay into the boundary. Then
\begin{equation}
X''+\left(\frac{\mu}{L}\right)^2X=0,\quad X'(0)=0
\;\Rightarrow\;
X=\cos\!\left(\frac{\mu x}{L}\right);
\qquad
Y''-\left(\frac{\mu}{L}\right)^2Y=0,\;Y\to0
\;\Rightarrow\;
Y=e^{-\mu y/L}.
\end{equation}
Hence
\begin{equation}
C_{\Ox}(x,y)=A\,\cos\!\left(\frac{\mu x}{L}\right)e^{-\mu y/L},
\label{eq:mode}
\end{equation}
where $A$ is an amplitude and $\mu$ is to be fixed by the interfacial condition.
Different values of $\mu$ give different modes; the fundamental mode has the smallest
$\mu$ and the slowest decay in $y$, and by A5 it is the one retained.

\subsection{Averages and flux of the mode}

From \eqref{eq:mode}, with the common factor $e^{-\mu y/L}$ absorbed into $A$,
\begin{align}
\Cb&=\frac{A}{L}\int_{0}^{L}\cos\!\left(\frac{\mu x}{L}\right)\dd x
    =A\,\frac{\sin\mu}{\mu},
\label{eq:mode-avg}\\
\Ci&=A\cos\mu,
\label{eq:mode-int}\\
\Cc&=A,
\label{eq:mode-center}\\
\Jr&=-\D\,\frac{\partial C_{\Ox}}{\partial x}\bigg|_{x=L}
    =\frac{\D\,A\,\mu\sin\mu}{L}.
\label{eq:mode-flux}
\end{align}
Eliminating the amplitude with \eqref{eq:mode-avg}, $A=\mu\Cb/\sin\mu$, gives the
two relations of \eqref{eq:result}:
\begin{equation}
\Jr=\frac{\D\,\mu^{2}\,\Cb}{L},
\qquad
\Ci=\Cb\,\mu\cot\mu,
\qquad
\Cc=\frac{\mu\,\Cb}{\sin\mu}.
\label{eq:mode-relations}
\end{equation}
Both $\Jr$ and $\Ci$ are expressed through $\Cb$ and $\mu$ only. Note that
$\mu\cot\mu$ decreases monotonically from $1$ at $\mu=0$ to $0$ at $\mu=\pi/2$, so
$0\le\Ci\le\Cb$ for $\mu\in[0,\pi/2]$.

\subsection{Consistency with the averaged equation}

Substituting the mode into the steady form of \eqref{eq:1d-working} with constant
$L$ and $\D$,
\begin{equation}
\D L\,\frac{\partial^2\Cb}{\partial y^2}
 =\D L\left(\frac{\mu}{L}\right)^2\Cb
 =\frac{\D\,\mu^2\,\Cb}{L}
 =\Jr ,
\end{equation}
with no residual. The averaged equation and the mode closure are therefore mutually
consistent; no additional approximation is introduced by combining them.

\section{The equation for $\mu$}
\label{sec:mu}

\subsection{General (nonlinear) reaction law}

Imposing the interfacial condition of \eqref{eq:laplace} on the mode, and using
\eqref{eq:mode-relations} to express both sides through $\Cb$ and $\mu$,
\begin{equation}
g(\mu)\equiv\frac{\D\,\mu^{2}\,\Cb}{L}-\Jr\bigl(\Cb\,\mu\cot\mu\bigr)=0 .
\label{eq:g}
\end{equation}
This is Eq.~\eqref{eq:mu-eq}. It is a single scalar equation in $\mu$; the inputs are
the local values of $\Cb$, $L$, $\D$ and the metal concentrations that enter
$\Jr(\cdot)$.

\paragraph{Existence and uniqueness.} On $(0,\pi/2)$ the first term increases
monotonically from $0$; the second term decreases monotonically, because $\Jr(\cdot)$ is
nondecreasing and $\mu\cot\mu$ is decreasing. Hence $g$ is strictly increasing, with
\begin{equation}
g(0^+)=-\Jr(\Cb),
\qquad
g(\pi/2)=\frac{\pi^2\D\,\Cb}{4L}-\Jr(0).
\end{equation}
If $\Jr(\Cb)>0$ (oxidation possible at the local average concentration) and
$\Jr(0)\le0$ (no consumption at zero oxygen, which holds for all four rate laws,
including the magnetite law below its threshold), there is exactly one root
$\mu^\ast\in(0,\pi/2)$. If $\Jr(\Cb)\le0$, no oxygen is consumed: set $\mu=0$,
$\Jr=0$, $\Ci=\Cb$.

\paragraph{Newton iteration.} The derivative is
\begin{equation}
g'(\mu)=\frac{2\D\,\mu\,\Cb}{L}
        -\Jr'\bigl(\Ci\bigr)\,\Cb\,\frac{\dd}{\dd\mu}\bigl(\mu\cot\mu\bigr),
\qquad
\frac{\dd}{\dd\mu}\bigl(\mu\cot\mu\bigr)=\frac{\sin\mu\cos\mu-\mu}{\sin^2\mu}<0,
\end{equation}
so $g'>0$ throughout and Newton's method converges from any bracketed start. A
safeguarded Newton (fall back to bisection if an iterate leaves $(0,\pi/2)$) handles
the kink of the magnetite law at $\Ci=C_{\Ox}^{\mathrm{th}}$. A good initial guess is
the linearized value of Sec.~\ref{sec:linear}.

\paragraph{Algorithm per grid point and time step.}
\begin{enumerate}
  \item Inputs: $\Cb$, $L$, $\D$, $C_{\mathrm{Cr}}^{\GB}$, $C_{\mathrm{Fe}}^{\GB}$,
        $C_{\mathrm{Si}}^{\GB}$.
  \item If $\Jr(\Cb)\le0$: $\mu=0$, $\Jr=0$, $\Ci=\Cb$. Stop.
  \item Solve $g(\mu)=0$ on $(0,\pi/2)$ by safeguarded Newton.
  \item $\Jr=\D\mu^2\Cb/L$; $\Ci=\Cb\,\mu\cot\mu$; the individual $\Jr^{k}(\Ci)$
        give the metal sinks, Eqs.~(\texttt{eq:sink-cr})--(\texttt{eq:sink-si}) of
        \texttt{main.tex}, and the growth of each $L_k$.
  \item Use $\Jr$ as the sink in \eqref{eq:1d-working}.
\end{enumerate}
This replaces the present quasi-steady closure, Eq.~(\texttt{eq:qss-closure}), which
also solves one scalar equation per point; the unknown is now $\mu$ instead of $\Ci$.

\subsection{Linearized reaction law}
\label{sec:linear}

For $\Jr=k\,\Ci$, Eq.~\eqref{eq:g} becomes
\begin{equation}
\frac{\D\mu^2}{L}=k\,\mu\cot\mu
\qquad\Longleftrightarrow\qquad
\mu\tan\mu=\frac{kL}{\D}\equiv\mathrm{Da},
\label{eq:mu-linear}
\end{equation}
the classical eigenvalue condition for a Robin boundary. $\mathrm{Da}$ is the ratio of
the lateral diffusion resistance $L/\D$ to the reaction resistance $1/k$. The
concentration $\Cb$ drops out: $\mu$ depends only on $\mathrm{Da}$.

\begin{table}[h]
\centering
\begin{tabular}{rrrr}
\toprule
$\mathrm{Da}=kL/\D$ & $\mu$ & $\Ci/\Cb=\mu\cot\mu$ & $\Jr L/(\D\Cb)=\mu^2$ \\
\midrule
$0$      & $0$     & $1$     & $0$ \\
$0.1$    & $0.311$ & $0.968$ & $0.097$ \\
$1$      & $0.860$ & $0.740$ & $0.740$ \\
$10$     & $1.429$ & $0.204$ & $2.042$ \\
$\infty$ & $\pi/2$ & $0$     & $\pi^2/4\approx2.467$ \\
\bottomrule
\end{tabular}
\caption{Fundamental mode for a linear reaction law.}
\label{tab:mu}
\end{table}

\subsection{Limits}

\begin{itemize}
  \item \textbf{Reaction-limited}, $\mathrm{Da}\ll1$: $\mu^2\simeq\mathrm{Da}$, so
        $\Jr\simeq k\,\Cb$ and $\Ci\simeq\Cb$. The interface sees the average
        concentration; the lateral profile is nearly uniform, with a parabolic
        deviation of relative size $O(\mathrm{Da})$. The decay length along the
        boundary is $L/\mu\simeq\sqrt{\D L/k}\gg L$.
  \item \textbf{Diffusion-limited}, $\mathrm{Da}\gg1$: $\mu\to\pi/2$, so
        $\Jr\to\pi^2\D\Cb/(4L)$ and $\Ci\to0$. Oxygen is consumed as soon as it
        reaches the interface; the lateral profile is a quarter cosine; the decay
        length along the boundary is $2L/\pi\approx0.64L$.
\end{itemize}
Between these limits $\mu$ interpolates smoothly. The closure never needs a
shape factor: if one insists on writing $\Jr=f\,\D(\Cb-\Ci)/L$, then
$f=\mu^2/(1-\mu\cot\mu)$, which runs from $3$ to $\pi^2/4$, but $f$ plays no role in
the algorithm.

\section{Range of validity}
\label{sec:validity}

\begin{enumerate}
  \item \textbf{Higher modes.} The interfacial condition admits a sequence of
        eigenvalues $\mu_0<\mu_1<\dots$, with $\mu_n\in(n\pi,\,n\pi+\pi/2)$ for a
        Robin condition. Each decays as $e^{-\mu_n y/L}$. The first higher mode decays
        at least as fast as $e^{-\pi y/L}$, i.e.\ over a distance $\lesssim L/\pi
        \approx0.3L$, compared with $\ge2L/\pi$ for the fundamental. Higher modes are
        excited where the boundary data change abruptly: the surface entry, the
        oxidation tip, and any position where $L$, $\D$ or the metal concentrations
        jump. Away from such places, only the fundamental survives (A5).
  \item \textbf{Slowly varying geometry (A3).} The mode is derived for a strip of
        constant width. For a wedge-shaped oxide the closure is applied with the local
        $L(y)$; the error is $O(\partial L/\partial y)$, the same order as the terms
        dropped in \eqref{eq:1d-exact}.
  \item \textbf{Quasi-steady lateral profile (A4).} Requires $L^2/\D$ to be short
        compared with the time scales of $\Cb(t)$ and $L(t)$. This can be checked
        from the simulation output.
  \item \textbf{Early stage.} When the oxidized zone is as deep as it is wide, the
        decay length is comparable to $L$ and the whole problem is two-dimensional.
        The one-dimensional description is meaningful only once the oxidized zone is
        slender; the initial aspect-ratio-one stage lies outside its scope.
  \item \textbf{Uniform $\D$ across the section (A1).} This is what allows the same
        $\D$ in \eqref{eq:1d-working} and \eqref{eq:laplace}. It follows from the
        local mixing of phases assumed in \texttt{main.tex}.
\end{enumerate}

\section{Summary of the model change}
\label{sec:summary}

Current formulation in \texttt{main.tex}:
\begin{align}
\frac{\partial C_{\Ox}}{\partial t}
 &=\frac{\partial}{\partial y}\!\left(\D\frac{\partial C_{\Ox}}{\partial y}\right)
   -\frac{\Jr}{\delta_{\GB}}, &
\Jr &=\D\,\frac{C_{\Ox}-\Ci}{\sum_k L_k}=\sum_k\Jr^{k}(\Ci).
\end{align}
Proposed formulation (cross-section average over the whole oxide, fundamental-mode
closure):
\begin{align}
\frac{\partial}{\partial t}\bigl(L\,\Cb\bigr)
 &=\frac{\partial}{\partial y}\!\left(\D\,L\,\frac{\partial\Cb}{\partial y}\right)-\Jr, &
\Jr &=\frac{\D\,\mu^{2}\,\Cb}{L},\quad
\frac{\D\,\mu^{2}\,\Cb}{L}=\sum_k\Jr^{k}\!\bigl(\Cb\,\mu\cot\mu\bigr).
\end{align}
What changes: (i) the storage width and along-boundary conductance are the
time-dependent half-width $L$ instead of the fixed $\delta_{\GB}$, with $L$ inside the
time derivative; (ii) the lateral closure is the fundamental Laplace mode, valid from
the reaction-limited to the diffusion-limited regime, with $\mu$ obtained pointwise
by Newton's method; (iii) $L$ is the mid-plane-to-interface distance. What does not
change: the reaction laws, the metal sinks, the definition of $\D$ as a local
volume-fraction average, and the structure of the pointwise scalar solve.

\end{document}
